\section{Results}

\subsection{Bulk Floc Diameter and Settling Velocity}

Before examining the fractal structure of individual flocs, we first report population-level trends in floc diameter and settling velocity across experimental conditions (Fig.~\ref{fig:basic}). For each experiment, we computed the median equivalent circular diameter and median settling velocity of all detected flocs, with error bars indicating the interquartile range (IQR).

Median floc diameter decreases monotonically with shear velocity, from 43.2~$\mu$m at $u_* = 0.02$~m\,s$^{-1}$ to 28.0~$\mu$m at $u_* = 0.03$~m\,s$^{-1}$ and 26.8~$\mu$m at $u_* = 0.04$~m\,s$^{-1}$ (Fig.~\ref{fig:basic}a), with the largest reduction occurring between the two lowest shear velocities. This trend is consistent with the equilibrium flocculation model of \citeA{nghiem2022mechanistic}, which predicts a positive linear scaling between floc diameter and the Kolmogorov microscale, $\eta$. Because $\eta$ decreases with increasing shear velocity, stronger turbulence is expected to limit floc growth through more efficient breakage of aggregates that approach the size of the smallest eddies. Floc diameter shows a weaker response to organic matter content, increasing slightly from 42.0~$\mu$m at 1~wt\% OM to 44.6~$\mu$m at 2~wt\% OM before declining to 38.8~$\mu$m at 3~wt\% OM (Fig.~\ref{fig:basic}b). The \citeA{nghiem2022mechanistic} model predicts that increasing organic matter should enhance bridging flocculation and increase floc diameter for the range of organic cover fraction explored here ($\theta < 0.5$), so the observed insensitivity or slight decrease at higher OM is not captured by their framework.

The settling velocity trends diverge more sharply from constant-fractal-dimension model predictions. \citeA{nghiem2022mechanistic} predicts that settling velocity should decrease monotonically with increasing shear velocity through the linear dependence on $\eta$. Instead, we observe a non-monotonic response: median settling velocity increases from 0.30~mm\,s$^{-1}$ at $u_* = 0.02$~m\,s$^{-1}$ to 0.39~mm\,s$^{-1}$ at $u_* = 0.03$~m\,s$^{-1}$ before declining to 0.32~mm\,s$^{-1}$ at $u_* = 0.04$~m\,s$^{-1}$ (Fig.~\ref{fig:basic}c). Floc diameter decreased over this same range, so the simultaneous increase in settling velocity between $u_* = 0.02$ and 0.03~m\,s$^{-1}$ implies that flocs became denser---a change not accounted for by a model with fixed fractal dimension. Similarly, \citeA{nghiem2022mechanistic} predicts that increasing organic matter should increase settling velocity through enhanced particle binding. We observe the opposite: settling velocity decreases monotonically from 0.78~mm\,s$^{-1}$ at 1~wt\% OM to 0.67~mm\,s$^{-1}$ at 2~wt\% OM and 0.52~mm\,s$^{-1}$ at 3~wt\% OM (Fig.~\ref{fig:basic}d).

These discrepancies point to a limitation of assuming a constant fractal dimension. In the fractal settling framework (Eq.~\eqref{eq:floc_settling}), settling velocity depends on both floc diameter and the fractal dimension through the effective density relation (Eq.~\eqref{eq:effective_density}): a decrease in $n_f$ produces more porous, lighter flocs that settle more slowly for a given diameter. If shear velocity simultaneously reduces floc diameter but increases settling velocity, the fractal dimension must be increasing (more compact flocs). Conversely, if organic matter reduces settling velocity while floc diameter remains relatively unchanged, the fractal dimension must be decreasing (more open, porous flocs). The bulk trends in Fig.~\ref{fig:basic} therefore motivate a detailed analysis of how fractal structure varies with shear velocity and organic matter---the subject of the following sections.

\begin{figure}[htbp]
    \centering
    \includegraphics[width=0.9\textwidth]{figures/flocpaperfigure_basic.pdf}
    \caption{Median floc diameter (blue, top row) and median settling velocity (red, bottom row) as functions of shear velocity (left column; Experiments~3--5, at 2~wt\% OM) and organic matter content (right column; Experiments~6--8, at $u_* = 0.04$~m\,s$^{-1}$). Error bars indicate interquartile ranges. Floc diameter decreases with increasing shear, consistent with turbulence-limited floc growth \cite{nghiem2022mechanistic}. However, the settling velocity trends---non-monotonic with shear and decreasing with organic matter---cannot be explained by a constant-fractal-dimension model and instead point to systematic variation in floc structure across conditions.}
    \label{fig:basic}
\end{figure}

\subsection{Stratified Inference of 3D Fractal Dimension}

To infer a three-dimensional fractal dimension from settling data, we group flocs into bins based on their two-dimensional box-counting fractal dimension, $n_f^{(2\mathrm{D})}$, measured from images. Flocs that share a similar projected geometry are expected to share a similar 3D structure, so the conditional $w_s$--$d_f$ slope within each bin isolates a narrower range of true $n_f$ values than a bulk regression over the entire population. The bin width in $n_f^{(2\mathrm{D})}$ acts as a coarse-graining scale rather than a physical control parameter; a sensitivity analysis varying the bin width from 0.025 to 0.1 confirms that finer stratification does not reveal qualitatively different settling regimes but instead inflates variance as bin populations shrink (see Supplementary Material). We therefore adopt a bin width of 0.1 for all subsequent analyses.

We applied this approach to all six floc datasets (Experiments~3--8), encompassing approximately 12\,500 individual floc measurements. Within each 0.1-wide $n_f^{(2\mathrm{D})}$ bin, outliers were removed with a standard interquartile range (IQR) filter on $\log_{10} w_s$: particles falling outside $[Q_1 - 1.5\,\mathrm{IQR},\; Q_3 + 1.5\,\mathrm{IQR}]$ were excluded. Bins retaining fewer than 50 particles were discarded, yielding five to eight usable bins per experiment. The resulting bins span $n_f^{(2\mathrm{D})} \approx 1.0$--$1.8$, with a median population of roughly 190 flocs per bin; the most populated bins ($n_f^{(2\mathrm{D})} \approx 1.4$--$1.6$) contain several hundred to over one thousand particles, while tail bins near $n_f^{(2\mathrm{D})} \lesssim 1.2$ or $\gtrsim 1.7$ are more sparsely sampled. Within each bin, a Bayesian MCMC regression in log--log space yields a posterior distribution of the slope $m$; the three-dimensional fractal dimension follows as $n_f = m + 1$.

The main panels of Fig.~\ref{fig:7} show the log--log scatter of settling velocity versus floc diameter for each experiment, color-coded by $n_f^{(2\mathrm{D})}$ bin. The stratified regression lines fan outward from a common convergence region near the smallest floc sizes: bins with higher $n_f^{(2\mathrm{D})}$ (lighter colors) yield systematically steeper slopes, corresponding to higher $n_f$. This pattern is consistent across all six experiments, confirming that the 2D box-counting dimension captures real structural differences that manifest in distinct settling behavior. The aggregated regression (orange dashed line) threads through the center of the data cloud but collapses this spread into a single slope. Inset panels plot the bin-median $n_f$ against the bin-midpoint $n_f^{(2\mathrm{D})}$, with 95\% credible intervals from the MCMC posteriors. Per-bin values of $n_f$ range from approximately 1.1--1.4 at low $n_f^{(2\mathrm{D})}$ to 1.9--2.4 at high $n_f^{(2\mathrm{D})}$, whereas the aggregated fit returns a single value near 1.7--2.0 for each experiment. A linear fit through the bin medians yields $R^2 = 0.81$--$0.88$ for the shear-velocity series (Figs.~\ref{fig:7}a--c) and $R^2 = 0.47$--$0.85$ for the organic-matter series (Figs.~\ref{fig:7}d--f); the lower bound corresponds to the 2\% OM experiment, where one interior bin departs from the otherwise monotonic trend.

\begin{figure}[htbp]
    \centering
    \includegraphics[width=1\textwidth]{figures/flocfigure7.pdf}
    \caption{Stratified and aggregated (bulk) settling-velocity fits for all floc datasets. (a--c)~Varying shear velocity ($u_* = 0.02$, 0.03, 0.04~m\,s$^{-1}$). (d--f)~Varying organic matter content (1\%, 2\%, 3\%). Data are colored by $n_f^{(2\mathrm{D})}$ bin; solid lines show per-bin MCMC median regressions and the orange dashed line is the aggregated fit. Insets: inferred $n_f$ versus bin-midpoint $n_f^{(2\mathrm{D})}$ with 95\% credible intervals; the orange band marks the aggregated $n_f$ and its 95\% posterior interval.}
    \label{fig:7}
\end{figure}

\subsection{Comparison of Stratified and Aggregated Fractal Dimensions}

To compare how the two approaches summarize structural variability, we construct a single representative $n_f$ for each experiment under each method (Fig.~\ref{fig:8}). In the aggregated approach, a single Bayesian regression is fit to all particles within an experiment; the reported median and error bars reflect the posterior distribution of that one slope parameter. In the stratified approach, each $n_f^{(2\mathrm{D})}$ bin yields its own posterior for $n_f$; we draw from each bin's posterior in proportion to the number of particles it contains and combine the draws into a single population-weighted distribution. The reported median and interquartile range (IQR) therefore reflect real particle-to-particle heterogeneity in fractal structure, rather than parameter uncertainty from a single fit.

Figure~\ref{fig:8}a shows the resulting $n_f$ as a function of shear velocity. The stratified median increases monotonically from 1.77 at $u_* = 0.02$~m\,s$^{-1}$ to 1.95 at 0.03~m\,s$^{-1}$ and 2.12 at 0.04~m\,s$^{-1}$, a total shift of $\Delta n_f \approx 0.35$ that is well described by the linear fit $n_f = 17.4\,u_* + 1.42$. In contrast, the aggregated median remains effectively flat ($n_f = 1.0\,u_* + 1.94$; $\Delta n_f \approx 0.02$), hovering near 2.0 regardless of forcing. The two approaches nearly coincide at intermediate shear but diverge at the extremes: at low shear the aggregated value exceeds the stratified median by 0.21, while at high shear it falls below by 0.12. This divergence arises because the bulk regression is dominated by the most populated $n_f^{(2\mathrm{D})}$ bins, which occupy a narrow band of slopes, whereas the stratified approach preserves the contributions of sparser low- and high-$n_f$ subpopulations.

Figure~\ref{fig:8}b shows an analogous comparison across organic matter (OM) content. Stratified medians decrease from 1.84 at 1\% OM to 1.66 at 2\% and 1.63 at 3\% ($n_f = -0.11\,\mathrm{OM} + 1.92$), a net decline of 0.21 over 1--3\% OM. The aggregated analysis again yields a flatter response ($n_f = -0.03\,\mathrm{OM} + 1.81$; $\Delta n_f \approx 0.05$), and at elevated OM the aggregated median exceeds the stratified value by up to 0.12 units, masking the shift toward more open, low-$n_f$ structures that organic matter promotes.

\subsection{Comparing 2D and 3D Fractal Dimension Relationships}

Figure~\ref{fig:8}c pools the per-bin data from all six experiments to examine the mapping between $n_f^{(2\mathrm{D})}$ and $n_f$. The 40 bin-median values (grey circles) are compared against three reference models. The idealized Euclidean relation $n_f = n_f^{(2\mathrm{D})} + 1$ (black dashed) assumes perfect space-filling and systematically overpredicts $n_f$ across the measured range. The empirical power-law (PL) and box-counting (BC) conversions of \citeA{wang2022fractal}, calibrated on computationally generated aggregates, bracket the data more closely: the BC model ($n_f = 0.8118\,(n_f^{(2\mathrm{D})})^{1.8054}$, solid red) tracks the upper envelope of the observations, while the PL model ($n_f = 0.2015\,(n_f^{(2\mathrm{D})})^{4.079}$, dashed red) underpredicts at moderate $n_f^{(2\mathrm{D})}$. Our data scatter between these two empirical curves, with the closest overall agreement to the BC conversion. This alignment is consistent with the fact that both our stratification variable and the Wang BC model are based on box-counting dimensions. Nevertheless, no single conversion curve captures the full spread, underscoring that the 2D--3D mapping is inherently system-specific and should be recalibrated when sediment composition or aggregation conditions change.

\begin{figure}[htbp]
    \centering
    \includegraphics[width=1\textwidth]{figures/flocpaperfigure8.pdf}
    \caption{Summary of stratified versus aggregated fractal dimensions and the 2D--3D relationship. (a)~Median $n_f$ versus shear velocity; (b)~median $n_f$ versus organic matter content. Red circles: stratified estimates (IQR error bars reflecting population heterogeneity); navy circles: aggregated estimates (error bars from the Bayesian posterior of a single bulk regression). Linear fits are annotated in matching colors. (c)~Per-bin $n_f$ (grey circles, pooled from all six experiments) versus $n_f^{(2\mathrm{D})}$, compared with the Euclidean relation $n_f = n_f^{(2\mathrm{D})} + 1$ (black dashed), the Wang et al.\ (2022) box-counting conversion (BC, red solid), and power-law conversion (PL, red dashed).}
    \label{fig:8}
\end{figure}

\subsection{Inferring Primary Particle Diameter ($d_p$)}

The stratified regressions also provide a means of estimating the effective primary particle diameter, $d_p$. In log--log space, Eq.~\eqref{eq:floc_settling} predicts that regression lines for different $n_f$ converge as $d_f \to d_p$: at the primary-particle scale, a floc reduces to a single grain whose settling velocity is independent of fractal structure. We exploit this convergence by computing all pairwise intersections of the per-bin regression lines from their MCMC posteriors. Repeating the calculation across thousands of posterior draws yields an ensemble of intersection diameters whose spread reflects both measurement scatter and finite-sample uncertainty.

Figure~\ref{fig:9}(a--c) compares the resulting $d_p$ posterior (red) with the known input grain-size distribution (navy) for the three shear-velocity experiments, and Fig.~\ref{fig:9}(e--g) shows the corresponding comparison for the three organic-matter experiments. In every case the inferred $d_p$ distribution is substantially narrower than the input, occupying roughly one order of magnitude versus nearly four for the source material. The $d_p$ posteriors peak in the coarse-silt range ($\sim$10--40~$\mu$m), overlapping with the second mode of the bimodal input distribution rather than the sub-micron clay mode. This suggests that the effective building blocks of the observed flocs are drawn predominantly from the silt fraction, while the finest clays either remain dispersed or contribute to floc mass without controlling the $w_s$--$d_f$ scaling.

Panels~(d) and (h) summarize the median $d_p$ (red, with IQR error bars) alongside the input-sediment median and IQR (grey) across the two parameter sweeps. With increasing shear velocity (panel~d), the inferred $d_p$ rises modestly from 21.7~$\mu$m at $u_* = 0.02$~m\,s$^{-1}$ to 23.0~$\mu$m at 0.03~m\,s$^{-1}$ and 29.8~$\mu$m at 0.04~m\,s$^{-1}$, consistent with stronger shear selectively destroying the most fragile aggregates and shifting the effective primary scale toward coarser grains. With increasing organic matter (panel~h), the median $d_p$ decreases monotonically from 34.2~$\mu$m at 1\% OM to 27.9~$\mu$m at 2\% and 17.6~$\mu$m at 3\%, indicating that organic binding enables finer particles to participate as effective building blocks. Across all conditions the inferred $d_p$ remains well below the input-sediment median of $\sim$39~$\mu$m and is confined to a much narrower range, reinforcing the interpretation that only a subset of the available grain sizes controls the floc settling relation.

\begin{figure}[htbp]
    \centering
    \includegraphics[width=1\textwidth]{figures/flocpaperfigure9.pdf}
    \caption{Inferred primary particle diameter ($d_p$) compared with the known input grain-size distribution. (a--c)~Shear-velocity series and (e--g)~organic-matter series: navy shading shows the input sediment (50:50 montmorillonite--silica mixture by volume); red shading shows the $d_p$ posterior from pairwise regression-line intersections; dashed vertical lines mark the respective medians. (d)~Median $d_p$ versus shear velocity; (h)~median $d_p$ versus organic matter content. Red circles with IQR error bars are the inferred $d_p$; grey circles with IQR error bars show the input-sediment distribution for reference.}
    \label{fig:9}
\end{figure}
